% appendix-b-glossary.tex — OpenBench Software Design Document. Written from .scratch/design-doc/brief.md on 2026-10-05.
% The domain vocabulary is generated by the Makefile (target glossary) from CONTEXT.md with the SRS generator; edit CONTEXT.md, not this file.
\section{Glossary}\label{app:glossary}

The domain vocabulary below is generated at build time from the project's shared vocabulary file and is identical to the glossary of the SRS; the design terms after it are specific to this document.

\begin{multicols}{2}
\footnotesize
\input{build/glossary-compact}

\glossgroup{Design terms}
\glossterm{Module}{A unit of the design with one area of responsibility: one Java package under \code{edu.cwru.openbench}, or the web client. A back-end module owns its tables and exposes only its application services and the DTOs they take and return.}{}
\glossterm{Application service}{A Spring bean that is a module's public entry point, such as \code{RequestLifecycleService}. Each method runs in one transaction unless marked read-only and carries the role check for its action.}{}
\glossterm{Port}{An interface declared by a lower module and implemented by a higher one, so the lower module can ask a question without depending upward. The only port is \code{ReductionGuard}, declared by Inventory and implemented by Allocation.}{}
\glossterm{Capability interface}{An optional interface an allocator implements beside \code{Allocator} when its Allocation Class has that behavior: \code{Waitlistable}, \code{Bumpable}, or \code{ApprovalDriven}. Requests asks the \code{AllocatorRegistry} whether a class has one.}{}
\glossterm{Allocator}{The component that evaluates, holds, commits, releases, and cancels lines of one Allocation Class behind the common \code{Allocator} interface: \code{ExclusiveAllocator}, \code{MeteredAllocator}, or \code{MatchedAllocator}.}{}
\glossterm{Hold (as a row state)}{The state \code{HELD} of an Exclusive occurrence or Metered allocation row. It blocks capacity exactly as \code{ALLOCATED} does but carries no Charge; committing turns it into \code{ALLOCATED}, and a release or lapse ends it.}{}
\glossterm{Occurrence}{One time window of an Exclusive line on one instance. A line with a weekly recurrence has several, each stored as its own row so that the exclusion constraint checks every window.}{}
\glossterm{Exclusion constraint}{A PostgreSQL constraint that rejects a row whose values conflict with an existing row under the given operators. It keeps Held or Allocated occurrences of one Resource from overlapping and Loans of one Person from overlapping.}{}
\glossterm{Type lock}{The row lock (\code{SELECT ... FOR UPDATE}) that every allocation decision takes on each Resource Type row it touches, in ascending id order, so that checks and writes for the same type never interleave.}{}
\glossterm{Transaction (database)}{A unit of database work that commits entirely or not at all; one per user or job action, at READ COMMITTED. Not to be confused with a Ledger Transaction, which is one row of the \code{transaction} table.}{}
\glossterm{Append-only table}{A table the application may insert into and read but never update or delete from, enforced by the database role's grants. The Ledger's \code{transaction} table is one; a correction is a new row.}{}
\glossterm{Migration}{A versioned SQL script that Flyway applies at start-up to bring the schema to the current version. Production and tests run the same migrations, and Hibernate only validates the result.}{}
\glossterm{Read model}{A read-only view of data shaped for display rather than for updates. The Dashboard module's SQL views are the only read model in the design.}{}
\glossterm{Scheduled job}{Work the server process runs on a timer rather than on a user action: \code{DueWorkJob} every 60 seconds and \code{PeriodRolloverJob} at each Period start and on start-up.}{}
\glossterm{ArchUnit test}{A test, written with the ArchUnit library, that inspects the compiled classes and fails the build when a dependency breaks the dependency rule, such as a lower module importing a higher one.}{}
\glossterm{Testcontainers}{A test library that starts a disposable PostgreSQL~16 container for the integration tests, so tests run against the same database and migrations as production.}{}
\glossterm{Problem detail}{The JSON error body of RFC~9457 (\code{application/problem+json}). Every error response carries one, with a stable \code{code} such as \code{INSUFFICIENT\_BALANCE}.}{}
\glossterm{Session cookie}{The HTTP cookie that identifies a signed-in Member's server-side session after form login. It is sent only to the same origin as the single-page application and expires after 12 hours idle.}{}
\glossterm{Single-page application}{A web client that loads once and then changes screens in the browser, fetching data from the REST API. Here, the React client served by the same server process.}{}
\end{multicols}
